\clearpage
\begingroup
\fancyhead[L]{\small THEORIE DER BEWEGUNG DES MONDES.}
\renewcommand{\headrulewidth}{0pt}
\par\noindent\hspace*{\parindent}Die zweite Gleichung wird
\[
\begin{aligned}
0={}&\frac{\mathrm{dd}u}{\mathrm{d}v\,\mathrm{d}u}
+\frac{2\mathrm{d}p}{p\,\mathrm{d}v}
-\tfrac32nn\frac{B^3}{R^3}\frac{\mathrm{d}u^2}{\mathrm{d}v^2}\Bigl\{\sin2E\Bigl[+2\varepsilon\sin(2E+a)\\
&\qquad-2\varepsilon\sin(2E-a)\Bigr]\Bigr\}.
\end{aligned}
\]
\noindent Statt \(\dfrac{B^3}{R^3}\) schreiben wir \(1-3\varepsilon\cos a\)
(wo wir das Apogäum der Sonne als unbeweglich ansehen dürfen, also
\(a=nv-\mathit{Const.}\)); und für \(\dfrac{\mathrm{d}u^2}{\mathrm{d}v^2}\),
\(1+4e\cos M\); folglich
\[
\begin{aligned}
0={}&\frac{\mathrm{dd}u}{\mathrm{d}v\,\mathrm{d}u}
+\frac{2\mathrm{d}p}{p\,\mathrm{d}v}-\tfrac32nn\Bigl\{\sin2E\\
&\qquad+2e\sin(2E+M)+2e\sin(2E-M)\\
&\qquad+\tfrac12\varepsilon\sin(2E+a)-\tfrac72\varepsilon\sin(2E-a)\Bigr\},
\end{aligned}
\]
\noindent folglich
\[
\begin{aligned}
\log\frac{pp\mathrm{d}u}{\mathrm{d}v}={}&\log A-(\tfrac34nn+\tfrac34n^3)\cos2E\\
&-nne\cos(2E+M)-(3nne+6n^3e)\cos(2E-M)\\
&\Bigl[-\tfrac38nn\varepsilon\cos(2E+a)
+\tfrac{21}{8}nn\varepsilon\cos(2E-a)\Bigr].
\end{aligned}
\]
\noindent Die Theile, die \(\varepsilon\) enthalten, lassen wir hier weg, weil diese
Coefficienten auch nach der Integration nur von der 3\textsuperscript{ten} Ordnung sind.
Hieraus also
\SourceLeaf{629}
\[
\begin{aligned}
\frac{pp\mathrm{d}u}{\mathrm{d}v}={}&A\Bigl\{1-\tfrac34nn\cos2E-3nne\cos(2E-M)\\
&-nne\cos(2E+M)\\
&\Bigl[\mathbin{\vcenter{\hbox{\rule{.33em}{.04em}\kern.025em\rule{.33em}{.04em}}}}
\tfrac38nn\varepsilon\cos(2E+a)
+\tfrac{21}{8}nn\varepsilon\cos(2E-a)\Bigr]\Bigr\}.
\end{aligned}
\]
\noindent Ferner wird
\[
\begin{aligned}
\psi={}&nn\Bigl(\tfrac12+\tfrac32\cos2E-\tfrac32\varepsilon\cos a\\
&\Bigl[+\tfrac34\varepsilon\cos(2E+a)-\tfrac{21}{4}\varepsilon\cos(2E-a)\Bigr]\Bigr)
\end{aligned}
\]
\noindent und statt \(p\), \(h(1-e\cos M)\) gesetzt,
\[
\begin{aligned}
\frac{\psi}{p^3}={}&\frac{nn}{h^3}\Bigl\{\tfrac12+\tfrac32e\cos M+\tfrac32\cos2E
-\tfrac32\varepsilon\cos a\\
&+\tfrac94e\cos(2E+M)+\tfrac94\cos(2E-M)\\
&\Bigl[+\tfrac34\varepsilon\cos(2E+a)-\tfrac{21}{4}\varepsilon\cos(2E-a)\Bigr]\Bigr\};
\end{aligned}
\]
\noindent endlich
\[
\pi=\tfrac32nn\sin2E
\]
\[
\frac{\mathrm{d}p}{p^4\,\mathrm{d}v}\pi
=\frac{nn}{h^3}\Bigl\{\tfrac34e\cos(2E-M)-\tfrac34e\cos(2E+M)\Bigr\}.
\]
\noindent Hieraus wird folglich die erste Gleichung:
\[
\begin{aligned}
0={}&\frac{\mathrm{dd}p}{\mathrm{d}v^2}+p-AA\Bigl\{1-\tfrac32nn\cos2E\\
&-6nne\cos(2E-M)-2nne\cos(2E+M)\\
&\Bigl[-\tfrac34nn\varepsilon\cos(2E+a)+\tfrac{21}{4}nn\varepsilon\cos(2E-a)\Bigr]\Bigr\}\\
&\times\Bigl\{k(1-\tfrac34gg+\tfrac34gg\cos2\lambda)\\
&-\frac{nn}{h^3}\Bigl(\tfrac12+\tfrac32e\cos M+\tfrac32\cos2E
-\tfrac32\varepsilon\cos a\\
&+3e\cos(2E+M)+\tfrac32e\cos(2E-M)\\
&\Bigl[+\tfrac34\varepsilon\cos(2E+a)-\tfrac{21}{4}\varepsilon\cos(2E-a)\Bigr)\Bigr]\Bigr\}.
\end{aligned}
\]
\noindent Schreibt man hier statt \(\dfrac1{h^3}\) den oben gefundenen Werth
\(k(1-3ee-\tfrac34gg)\) und vernachlässigt die höhern Potenzen, so wird endlich
\[
\begin{aligned}
0={}&\frac{\mathrm{dd}p}{\mathrm{d}v^2}+p-AAk(1-\tfrac12nn-\tfrac34gg)\\
&\Bigl\{1+\tfrac34gg\cos2\lambda-\tfrac32nne\cos M-3nn\cos2E\\
&+\tfrac32nn\varepsilon\cos a-\tfrac{15}{2}nne\cos(2E-M)-5nne\cos(2E+M)\\
&\Bigl[-\tfrac32nn\varepsilon\cos(2E+a)+\tfrac{21}{2}nn\varepsilon\cos(2E-a)\Bigr]\Bigr\}.
\end{aligned}
\]
\noindent Hier zeigt nun die Integration sogleich eine progressive Bewegung des
Mondsapogäum, da ohne diese \(\dfrac{\mathrm{dd}p}{\mathrm{d}v^2}+p\)
das Glied \(\cos M\) nicht enthalten könnte. Setzt man nemlich
\(\dfrac{\mathrm{d}M}{\mathrm{d}v}=\eta\), so wird die Integration geben
\[
\begin{aligned}
p={}&AAk(1-\tfrac12nn-\tfrac34gg)\Bigl\{1-\tfrac14gg\cos2\lambda
+nn\cos2E+\tfrac32nn\varepsilon\cos a\\
&-\frac{15}{2}nne\frac{\cos(2E-M)}{1-(2-2n-\eta)^2}\\
&+\tfrac58nne\cos(2E+M)-\tfrac32nne\frac{\cos M}{1-\eta\eta}\\
&\Bigl[+\tfrac12nn\varepsilon\cos(2E+a)-\tfrac72nn\varepsilon\cos(2E-a)\Bigr]\Bigr\}\\
&+\mathit{Const.}\cos(v+\mathit{Const.}).
\end{aligned}
\]
\clearpage
\endgroup
